\chapter{Design Process}\label{ch:design-process}
\section{Getting Started}
Collect the information that will shape the controls work as it becomes available: project goals and schedule, P\&IDs, equipment and valve lists, process descriptions, layout drawings, utility and power information, area classification, and network requirements. Keep a list of open questions, assumptions, and who will answer them. The bulk of controls engineering starts near the end of Stage 3; Stage 2 is for early discussions and planning.

\section{Project Execution Phases}
Use the project execution phase numbers and names below when planning controls work. These are practical handoffs, so work may overlap. Each stage starts with a short summary, then lists the \textbf{required inputs}, the \textbf{work performed}, and the \textbf{outputs}. Check the inputs before promising the outputs. Stage 5A is an added controls handoff between Construction and Commissioning; it does not renumber the company phases.

% Local helpers for the stage lists
\newcommand{\stagelist}[1]{\par\smallskip\noindent\textbf{#1}\par\vspace{-2pt}}

\subsection{Stage 0 --- Project Evaluation}
\textit{The business is deciding whether the project is worth pursuing. Controls involvement is light: if asked, give a quick, high-level opinion on how the system could be automated and flag any obvious risks. There is normally no detailed controls design in this stage.}

\stagelist{Required Controls Inputs}
\begin{enumerate}[itemsep=1pt]
  \item Proposed process
  \item Rough scope
  \item Likely site
  \item Any early questions about automation
\end{enumerate}

\stagelist{Work Performed}
\begin{enumerate}[itemsep=1pt]
  \item If asked, give a high-level opinion on the control platform and how the system would integrate with existing equipment.
  \item If asked, give a high-level opinion on panel scope.
  \item Point out any obvious technical or schedule risks.
\end{enumerate}

\stagelist{Controls Outputs}
\begin{enumerate}[itemsep=1pt]
  \item High-level opinion on platform and integration approach (if requested)
  \item High-level opinion on panel scope (if requested)
  \item Flagged technical or schedule risks (if requested)
\end{enumerate}

\subsection{Stage 1 --- Project Initiation}
\textit{The project has been approved and the team is forming. Controls gets read into the project, learns who owns which documents, and starts a list of open questions.}

\stagelist{Required Controls Inputs}
\begin{enumerate}[itemsep=1pt]
  \item Approved project direction
  \item Team contacts
  \item Overall project goal
  \item General schedule
  \item What is known about the equipment and site
\end{enumerate}

\stagelist{Work Performed}
\begin{enumerate}[itemsep=1pt]
  \item Get read into the project: goals, schedule, site, and equipment.
  \item Record open controls questions and key interfaces (other systems, vendors, utilities).
  \item Confirm who owns the P\&IDs, Control Narrative, Interlock Matrix, and layout drawing.
  \item Identify the information needed for preliminary engineering.
\end{enumerate}

\stagelist{Controls Outputs}
\begin{enumerate}[itemsep=1pt]
  \item List of open controls questions and key interfaces
  \item Confirmed owner of the P\&IDs
  \item Confirmed owner of the Control Narrative
  \item Confirmed owner of the Interlock Matrix
  \item Confirmed owner of the layout drawing
\end{enumerate}

\subsection{Stage 2 --- Preliminary Engineering}
\textit{The process design is taking shape. Controls works with the process engineer on the broad operating approach and sketches possible panel, platform, and network arrangements. These are high-level discussions, not the start of full electrical design or PLC programming.}

\stagelist{Required Controls Inputs}
\begin{enumerate}[itemsep=1pt]
  \item Rough P\&IDs
  \item Early process description
  \item Likely system location
  \item Preliminary schedule
\end{enumerate}

\stagelist{Work Performed}
\begin{enumerate}[itemsep=1pt]
  \item Discuss the broad operating approach with the process engineer, with assumptions clearly marked.
  \item Sketch possible panel locations, control platform, and network arrangement.
  \item Start the FAT Checkout document outline; specific checks are added as the design develops.
\end{enumerate}

\stagelist{Controls Outputs}
\begin{enumerate}[itemsep=1pt]
  \item Sketches of possible panels, control platform, and network arrangement (assumptions marked)
  \item FAT Checkout document outline (started; specific checks added as design develops)
\end{enumerate}

\subsection{Stage 3 --- Basic Engineering}
\textit{The P\&IDs and layout are firm enough to count I/O and pick a panel concept. This is when the bulk of controls engineering begins, and when long-lead parts (often PLC hardware) are ordered once the project team approves them.}

\stagelist{Required Controls Inputs}
\begin{enumerate}[itemsep=1pt]
  \item Developing P\&IDs
  \item Preliminary system layout
  \item Equipment and utility information
  \item Enough scope to identify likely long-lead hardware
\end{enumerate}

\stagelist{Work Performed}
\begin{enumerate}[itemsep=1pt]
  \item Build a preliminary I/O count and list from the P\&IDs and review it with the project team.
  \item Near the end of the stage, decide whether one central panel or remote panels make sense.
  \item Identify how 240/480 VAC power equipment will be separated from 120 VAC and 24 VDC controls.
  \item Once the project team approves the quantities and assumptions, order long-lead items (often PLC components).
  \item Record spare capacity and anything that is still uncertain. Detailed schematics and the final I/O list belong in Stage 4.
\end{enumerate}

\stagelist{Controls Outputs}
\begin{enumerate}[itemsep=1pt]
  \item Preliminary I/O count and panel concept
  \item Preliminary I/O list, reviewed with the project team
  \item Long-lead item orders placed (often PLC components), once approved
  \item Record of spare capacity or remaining uncertainty
\end{enumerate}

\subsection{Stage 4 --- Detailed Engineering}\label{sec:stage4}
\textit{This is the main design stage. Using the finished P\&IDs, Control Narrative, Interlock Matrix, and layout, controls produces the full electrical package: I/O list, schematics, panel layouts, bill of materials, and network map.}

\stagelist{Required Controls Inputs}
\begin{enumerate}[itemsep=1pt]
  \item Completed P\&IDs
  \item Interlock Matrix (Table of Interlocks and Permissives)
  \item Control Narrative
  \item System layout drawing
  \item Applicable PHA results and approved controls recommendations (obtain before treating the interlock table as final)
\end{enumerate}

\stagelist{Work Performed}
\begin{enumerate}[itemsep=1pt]
  \item If an input is still in progress, draft what you can and keep its assumptions open for review.
  \item Use the physical layout to place panels and remote I/O (see Section~\ref{sec:local-remote-io}).
  \item Use the I/O list to size racks, cards, terminals, and power.
  \item Draw the schematics and panel layouts, and build the electrical bill of materials.
  \item Confirm high-voltage separation and the panel build package with the project team.
  \item Draw the network map and a preliminary cable schedule.
  \item Mock up the SCADA layout and operator navigation.
  \item \textit{\textcolor{futureblue}{[Future: write a protective-function design record for each trip (see Section~\ref{sec:protective-function-reqs}). This would be developed here alongside the Interlock Matrix and checked during Stage 6 commissioning.]}}
\end{enumerate}

\stagelist{Controls Outputs}
\begin{enumerate}[itemsep=1pt]
  \item Approved I/O list
  \item Electrical bill of materials
  \item Electrical schematics
  \item Panel layouts
  \item Network map
  \item Preliminary cable schedule
  \item Preliminary SCADA layout and operator navigation mockup (Visio, \href{https://www.drawio.com/}{draw.io}, or similar)
\end{enumerate}

\noindent\textit{Suggested handoff between Stages 4 and 5: Building (Pre-Construction).} The project team could use this name for the period when design is released, the panel shop is lined up, and material orders are completed. It is a proposed planning note, not an additional numbered company phase.

\subsection{Stage 5 --- Development \& Construction}
\textit{The design is released and things get built. The panel shop builds the panels while controls writes the PLC program and SCADA/HMI application in parallel. Late in the stage, the electrician connects the field devices to the panels.}

\stagelist{Required Controls Inputs}
\begin{enumerate}[itemsep=1pt]
  \item Released Stage 4 electrical package
  \item Approved purchasing and panel-shop scope
  \item Current Control Narrative and Interlock Matrix
  \item Access to the project development environment
  \item Construction schedule coordinated with mechanical and electrical contractors
\end{enumerate}

\stagelist{Work Performed}
\begin{enumerate}[itemsep=1pt]
  \item Receive the parts and have the contractor or panel shop build the electrical panels.
  \item Develop the PLC program and SCADA/HMI application in parallel with the physical build.
  \item Keep expanding the FAT Checkout document started in Stage 2. Its procedures and expected results should be complete by late Stage 5.
  \item Late in the stage, once the panels and mechanical system are built, the electrical contractor connects the equipment, sensors, motors, heaters, and other field devices to their panels. The electrician is responsible for correct installation and for verifying those connections (identification and continuity checks).
  \item Support the electrician's checkout, answer design questions, and record approved changes in the drawings and I/O list. The contractor's wiring checks do not replace controls I/O and functional testing.
\end{enumerate}

\stagelist{Controls Outputs}
\begin{enumerate}[itemsep=1pt]
  \item Built electrical panels
  \item PLC program, targeting about 95\% completion by end of Stage 5
  \item SCADA/HMI application, targeting about 95\% completion by end of Stage 5
  \item Finished network diagram
  \item FAT Checkout document with complete procedures and expected results
\end{enumerate}

\subsection{Stage 5A --- PLC and SCADA Completion}
\textit{A controls-only handoff stage. Controls finishes the PLC and SCADA applications, runs as much of the FAT Checkout as possible before going to the field, and makes clean, versioned backups.}

\stagelist{Required Controls Inputs}
\begin{enumerate}[itemsep=1pt]
  \item Near-complete Stage 5 applications
  \item Built panels
  \item Current field wiring and network information
  \item Completed FAT Checkout procedure
\end{enumerate}

\stagelist{Work Performed}
\begin{enumerate}[itemsep=1pt]
  \item Finish the PLC and SCADA/HMI applications and resolve development punch-list items.
  \item Make versioned backups.
  \item Run the available panel and integrated tests using the FAT Checkout document, and record actual results and open items.
  \item Clearly list any tests that can only be done in the field so they are picked up during commissioning.
\end{enumerate}

\stagelist{Controls Outputs}
\begin{enumerate}[itemsep=1pt]
  \item Completed PLC and SCADA/HMI programs
  \item Versioned backups
  \item FAT Checkout actual results and open items
  \item As-built connection information available at handoff
  \item List of remaining field-dependent tests, clearly identified for commissioning
\end{enumerate}

The completed PLC and SCADA programs, together with the as-built connection information available at handoff, are required inputs to Stage 6.

\subsection{Stage 6 --- Commissioning}
\textit{Everything is installed and wired. Controls checks every field signal end to end, then tests sequences, alarms, and interlocks with the project team before equipment is run.}

\stagelist{Required Controls Inputs}
\begin{enumerate}[itemsep=1pt]
  \item Completed PLC and SCADA/HMI programs from Stage 5A
  \item Versioned backups
  \item Built panels
  \item Installed mechanical equipment
  \item Verified field wiring
  \item Current I/O list and drawings
  \item FAT Checkout document with available results and open items
\end{enumerate}

\stagelist{Work Performed}
\begin{enumerate}[itemsep=1pt]
  \item Check each field signal through the PLC and SCADA: scaling, signal direction, valve action, motor rotation, fail states, communications, and safety circuits.
  \item Test sequences, alarms, permissives, interlocks, shutdowns, power recovery, and abnormal conditions with the project team.
  \item Run equipment only after wiring and safety checks are complete.
  \item Track punch-list items, log software changes, and make new versioned backups.
\end{enumerate}

\stagelist{Controls Outputs}
\begin{enumerate}[itemsep=1pt]
  \item Field signal check results (scaling, direction, valve/motor action, fail states, comms, safety circuits)
  \item Sequence, alarm, permissive, interlock, shutdown, and abnormal-condition test results
  \item Punch-list ownership and status
  \item Software change log
  \item New versioned backups
\end{enumerate}

\subsection{Stage 7 --- Startup/Qualification}
\textit{The process is brought online. Controls tunes loops, confirms the system behaves correctly under real operating conditions, supports qualification testing, and trains operators.}

\stagelist{Required Controls Inputs}
\begin{enumerate}[itemsep=1pt]
  \item Commissioning records
  \item Resolved critical punch-list items
  \item Approved operating limits
  \item Agreed startup plan
\end{enumerate}

\stagelist{Work Performed}
\begin{enumerate}[itemsep=1pt]
  \item Bring equipment, and then the process, online with the process engineer, operators, and vendors.
  \item Tune control loops.
  \item Confirm expected normal and fault behavior under operating conditions.
  \item Support the qualification tests required by the project.
  \item Train operators on the actual screens, alarms, and recovery steps.
\end{enumerate}

\stagelist{Controls Outputs}
\begin{enumerate}[itemsep=1pt]
  \item Tuned control loops
  \item Confirmed normal and fault behavior under operating conditions
  \item Qualification test support records
  \item Software revision log and list of remaining issues
  \item Trained operators
\end{enumerate}

\subsection{Stage 8 --- Project Closeout}
\textit{The project wraps up. Controls closes or hands off remaining items and stores a complete, final set of files so the next person can support the system.}

\stagelist{Required Controls Inputs}
\begin{enumerate}[itemsep=1pt]
  \item Startup and qualification results
  \item Field changes
  \item Open-item status
  \item Final panel and software configurations
\end{enumerate}

\stagelist{Work Performed}
\begin{enumerate}[itemsep=1pt]
  \item Close or formally transfer remaining controls items.
  \item Update drawings, lists, and documents to as-built.
  \item Store the approved set in the project location so the next person can support the system.
\end{enumerate}

\stagelist{Controls Outputs}
\begin{enumerate}[itemsep=1pt]
  \item Final PLC and SCADA source code and backups
  \item As-built schematics and panel layouts
  \item I/O and cable lists
  \item Final network diagram and address list
  \item Current Control Narrative and Interlock Matrix
  \item Test records
  \item Manuals
  \item Relevant configuration files
\end{enumerate}

\section{Typical Design Deliverables}
The phase outputs above are the working checklist. Scale the detail to the project, but keep the I/O list, Control Narrative, Interlock Matrix, electrical package, network diagram, software backups, checkout records, and as-built documentation in step with the actual system. Larger or higher-risk projects may also need formal calculations, cybersecurity records, alarm documentation, and validation records.

\section{Reviews and Changes}
Bring in the people who will design, build, operate, maintain, and support the system when their input is useful. Track important review comments and open items. When a late change affects hardware, safety, software behavior, testing, or documentation, check the impact and update the affected files.